\documentclass[a4paper,12pt]{article}
	\usepackage[ansinew]{inputenc}
	\usepackage{ae}	
	\usepackage[brazil]{babel}
	\usepackage{xcolor}
	\usepackage{graphicx}
		
	\graphicspath{
		{Figuras/EPS/}
		{Figuras/JPG/}
		{Figuras/PDF/}
		{Figuras/PNG/}
		}
		
	% O comando \input insere um arquivo .tex no seu documento. Neste caso o LaTeX SUBSTITUI a instrução abaixo pelo conteúdo do arquivo minhasconfiguracoes.tex. Aproveite para observar que este arquivo não tem preâmbulo nem o ambiente document. Por quê?
	\newcommand{\comando}[1]{{\normalfont\textbackslash\color{blue!50!black}#1}}
\newcommand{\pacote}[1]{{\normalfont\sffamily\color{orange!50!black}#1}}
\newcommand{\ambiente}[1]{{\normalfont\sffamily\color{green!50!black}#1}}
\newcommand{\contador}[1]{{\normalfont\sffamily\color{red!50!black}#1}}
\newcommand{\foreign}[1]{\textsl{#1}}
\newcommand\linhabase{\noindent\makebox[0pt][l]{\color{red}\rule{\textwidth}{0.1pt}}}
	
\begin{document}

	\begin{figure}
		\begin{minipage}[b]{0.4\textwidth}
			\includegraphics[width=\textwidth]{Tei-Gi}
		\end{minipage}\hfill
		\begin{minipage}[b]{0.5\textwidth}
			\caption{legenda da figura, colocada ao lado dela através de duas \ambiente{minipage}s: a primeira contém a figura, a segunda contém esta legenda e ambas compõem uma linha de texto... caixa, espaço elástico (\comando{hfill}) e caixa.}
			\label{fig:minipage}
		\end{minipage}
	\end{figure}

	Dentro de um ambiente \ambiente{minipage} (fig.~\ref{fig:minipage}), o comprimento \comando{textwidth} assume a largura da minipágina. Por isso no exemplo acima \comando{textwidth} define a largura da figura, de modo a ocupar toda a extensão horizontal da minipágina. O mesmo não acontece com o comando \comando{parbox} (fig.~\ref{fig:parbox}).

	\begin{figure}
		\parbox[b]{0.4\textwidth}{\includegraphics[width=0.4\textwidth]{Tei-Gi}}%
		\hfill
		\parbox[b]{0.5\textwidth}{\caption{legenda da figura, colocada ao lado dela através de duas \ambiente{minipage}s: a primeira contém a figura, a segunda contém esta legenda e ambas compõem uma linha de texto... caixa, espaço elástico (\comando{hfill}) e caixa.}
		\label{fig:parbox}}
	\end{figure}

\end{document}
